\documentclass[11pt,a4paper,oneside]{book}
\usepackage{90-faelles/latex/studiearkiv}
\hypersetup{pdftitle={Matematik – Bind 5: Geometri og tensorer}}
\begin{document}
\frontmatter
\begin{titlepage}
\centering
\vspace*{30mm}
{\sffamily\color{arkivblaa}\Large MATEMATIK TIL UNIVERSITETSFYSIK\par}
\vspace{16mm}
{\Huge\bfseries Bind 5\\[6mm]Geometri og tensorer\par}
\vspace{13mm}
{\Large Formler, begreber og matematisk sammenhæng\par}
\vfill
\begin{minipage}{0.8\textwidth}
Geometri i 2D og 3D, vektorfelter, tensorer og krumme rum.

Dette bind genbruger de samme LaTeX-kapitler som det samlede matematikkompendium. Planlagte kapitler er markeret som ufærdige. Der indgår ikke separate gennemregnede opgaveeksempler.
\end{minipage}
\vspace{14mm}
{\small Fysik-studiearkiv · Arbejdsudgave\par}
\end{titlepage}
\chapter*{Læsevejledning}
\addcontentsline{toc}{chapter}{Læsevejledning}
Dette bind er en emneopdelt udgave af det samlede matematikkompendium. Den generelle kilde og alle figurer vedligeholdes centralt under \texttt{00-matematik/chapters} og \texttt{00-matematik/figures}. For hele overblikket samt henvisninger mellem bindene anvendes den samlede udgave. Ufuldstændige kapitler er tydeligt markerede.
\begingroup\small\tableofcontents\endgroup
\mainmatter
\chapter{Analytisk geometri i 2D og 3D}
\planlagt{Definitioner, notation, forudsætninger, centrale resultater og fysikforbindelser udvikles samlet.}
\section{Linjer, planer og afstande}
\leavevmode\par
% Redaktionelt: tilføj definitioner, gyldighed, relationer og relevante figurer.

\section{Keglesnit og kvadrikker}
\leavevmode\par
% Redaktionelt: tilføj definitioner, gyldighed, relationer og relevante figurer.

\section{Parametriserede kurver, krumning og torsion}
\leavevmode\par
% Redaktionelt: tilføj definitioner, gyldighed, relationer og relevante figurer.

\section{Koordinatsystemer og orientering}
\leavevmode\par
% Redaktionelt: tilføj definitioner, gyldighed, relationer og relevante figurer.

\chapter{Vektoranalyse}
\planlagt{Definitioner, notation, forudsætninger, centrale resultater og fysikforbindelser udvikles samlet.}
\section{Gradient, divergens og rotation}
\leavevmode\par
% Redaktionelt: tilføj definitioner, gyldighed, relationer og relevante figurer.

\section{Kurve- og fladeintegraler med orientering}
\leavevmode\par
% Redaktionelt: tilføj definitioner, gyldighed, relationer og relevante figurer.

\section{Gauss og Stokes med forudsætninger}
\leavevmode\par
% Redaktionelt: tilføj definitioner, gyldighed, relationer og relevante figurer.

\section{Ortogonal krum koordinatnotation}
\leavevmode\par
% Redaktionelt: tilføj definitioner, gyldighed, relationer og relevante figurer.

\chapter{Tensorregning}
\planlagt{Definitioner, notation, forudsætninger, centrale resultater og fysikforbindelser udvikles samlet.}
\section{Multilineære afbildninger og dualrum}
\leavevmode\par
% Redaktionelt: tilføj definitioner, gyldighed, relationer og relevante figurer.

\section{Indekser og transformationslove}
\leavevmode\par
% Redaktionelt: tilføj definitioner, gyldighed, relationer og relevante figurer.

\section{Kontraktion, symmetri og tensorprodukter}
\label{sec:tensor-kontraktion}
Kontraktion forbinder en tensor med et objekt af lavere orden ved at summere over ét øvre og ét nedre indeks. Den lettest genkendelige model er sporet af en lineær operator fra bind 1, hvor $A^i{}_i=\operatorname{tr}A$. Det er ikke det samme som at diagonaliserere en matrix: kontraktion summerer passende komponenter, hvorimod diagonalisering skifter basis, når det er muligt.

For en tensor af typen $(2,1)$ giver kontraktion eksempelvis
\[
T^{ij}{}_i=\sum_i T^{ij}{}_i,
\]
hvor resultatet har ét frit øvre indeks $j$ og derfor repræsenterer en vektor. Ved et skift af koordinater er det afgørende, at man parrer et kontravariant indeks med et kovariant indeks. At sætte to øvre indekser lig hinanden og summere giver ikke generelt en tensoriel operation.

Har man en ikke-degenereret metrik $g_{ij}$ med invers $g^{ij}$, kan man først sænke eller hæve indeks. For en kovariant tensor $S_{ij}$ fremkommer den metriske kontraktion
\[
g^{ij}S_{ij}.
\]
Dette bliver blandt andet væsentligt ved $R=g^{ij}R_{ij}$ i differentialgeometri. En fuld behandling af koordinattransformationer, symmetrisering og anti-symmetrisering er fortsat planlagt.

\paragraph{Tilbageblik.} Se kapitlet om lineær algebra i bind 1, afsnittet om sporet og diagonalisering, for den konkrete matrixfortolkning af kontraktion.

\section{Metrik, indeksflytning og orientering}
\leavevmode\par
% Redaktionelt: tilføj definitioner, gyldighed, relationer og relevante figurer.

\chapter{Differentialgeometri og matematik til GR}
\planlagt{Definitioner, notation, forudsætninger, centrale resultater og fysikforbindelser udvikles samlet.}
\section{Mangfoldigheder og koordinatkort}
\leavevmode\par
% Redaktionelt: tilføj definitioner, gyldighed, relationer og relevante figurer.

\section{Tangentrum og dualrum}
\leavevmode\par
% Redaktionelt: tilføj definitioner, gyldighed, relationer og relevante figurer.

\section{Tensorer og koordinattransformationer}
\leavevmode\par
% Redaktionelt: tilføj definitioner, gyldighed, relationer og relevante figurer.

\section{Metrikker og linjeelementer}
\leavevmode\par
% Redaktionelt: tilføj definitioner, gyldighed, relationer og relevante figurer.

\section{Kovariant differentiation og Christoffelsymboler}
\leavevmode\par
% Redaktionelt: tilføj definitioner, gyldighed, relationer og relevante figurer.

\section{Paralleltransport og holonomi}
\leavevmode\par
% Redaktionelt: tilføj definitioner, gyldighed, relationer og relevante figurer.

\section{Geodæter og affin parameter}
\leavevmode\par
% Redaktionelt: tilføj definitioner, gyldighed, relationer og relevante figurer.

\section{Riemanntensor og fortegnskonvention}
\leavevmode\par
% Redaktionelt: tilføj definitioner, gyldighed, relationer og relevante figurer.

\section{Riccitensor og skalarkrumning}
\leavevmode\par
% Redaktionelt: tilføj definitioner, gyldighed, relationer og relevante figurer.

\section{Geodætisk deviation}
\leavevmode\par
% Redaktionelt: tilføj definitioner, gyldighed, relationer og relevante figurer.

\section{Killingvektorer og Lie-afledte}
\leavevmode\par
% Redaktionelt: tilføj definitioner, gyldighed, relationer og relevante figurer.

\section{Differentialformer og Hodge-dualitet}
\leavevmode\par
% Redaktionelt: tilføj definitioner, gyldighed, relationer og relevante figurer.

\section{Bianchi-identiteterne}
\leavevmode\par
% Redaktionelt: tilføj definitioner, gyldighed, relationer og relevante figurer.

\section{Lorentzgeometri og introduktion til GR}
\leavevmode\par
% Redaktionelt: tilføj definitioner, gyldighed, relationer og relevante figurer.

\end{document}
